\chapter{テンソル計算}\label{tensor}

Egisonにはテンソル計算を便利に記述するために開発された機能が実装されている．
本章はこれらの機能を紹介していく．

\section{テンソルとは}

テンソルとは，ベクトルや行列を一般化したものである．
ベクトルは一階のテンソル，行列は二階のテンソルである．
ベクトルは一次元配列，行列は二次元配列としてプログラムでは表現されることが多い．
$n$階のテンソルは，$n$次元配列により表現される．

$
\begin{pmatrix}
a_{1} & a_{2} & a_{3}\\
\end{pmatrix}
$

$
\begin{pmatrix}
a_{11} & a_{12} & a_{13} \\
a_{21} & a_{22} & a_{23}
\end{pmatrix}
$

$
\begin{pmatrix}

\begin{pmatrix}
a_{111} & a_{112} & a_{113} \\
a_{121} & a_{122} & a_{123} 
\end{pmatrix}

\begin{pmatrix}
a_{211} & a_{212} & a_{213} \\
a_{221} & a_{222} & a_{223} 
\end{pmatrix}

\begin{pmatrix}
a_{311} & a_{312} & a_{313} \\
a_{321} & a_{322} & a_{323} 
\end{pmatrix}

\end{pmatrix}
$

\section{テンソルの添字記法}

テンソル計算は，ベクトルについてのテンソル積・アダマール積・内積の組み合わせでほとんど表現できる．
テンソル積は，ベクトルのすべての成分の組み合わせからなる行列を返す．
アダマール積は，対応する成分同士のみからなるベクトルを返す．
内積は，アダマール積の結果の成分を足し合わせた結果のスカラー値を返す．

$
\begin{pmatrix}
a_{1} & a_{2}
\end{pmatrix}
\otimes
\begin{pmatrix}
b_{1} & b_{2}
\end{pmatrix}
=
\begin{pmatrix}
a_{1} b_{1} & a_{1} b_{2} \\
a_{2} b_{1} & a_{2} b_{2}
\end{pmatrix}
$

$
\begin{pmatrix}
a_{1} & a_{2}
\end{pmatrix}
\circ
\begin{pmatrix}
b_{1} & b_{2}
\end{pmatrix}
=
\begin{pmatrix}
a_{1} b_{1} & a_{2} b_{2}
\end{pmatrix}
$

$
\begin{pmatrix}
a_{1} & a_{2}
\end{pmatrix}
\cdot
\begin{pmatrix}
b_{1} & b_{2}
\end{pmatrix}
=
a_{1} b_{1} + a_{2} b_{2}
$


たとえば，行列同士の掛け算は，Aの行とBの列についてはテンソル積，Aの列とBの行については内積をとる演算と解釈できる．

$
\begin{pmatrix}
a_{11} & a_{12} \\
a_{21} & a_{22}
\end{pmatrix}
\begin{pmatrix}
b_{11} & b_{12} \\
b_{21} & b_{22}
\end{pmatrix}
=
\begin{pmatrix}
a_{11} b_{11} + a_{12} b_{21} & a_{11} b_{12} + a_{12} b_{22} \\
a_{21} b_{11} + a_{22} b_{21} & a_{21} b_{12} + a_{22} b_{22}
\end{pmatrix}
$

ベクトルについては，掛け算の方法は，テンソル積・アダマール積・内積の3つだけだが，行列や行列よりさらに高階なテンソルとなると掛け算の方法が無数にある．
これらの掛け算のすべてに名前をつけるのは大変である．
そのため，数学者はシンボリックな添字をテンソルに付加し，その添字をパラメーターにしてこれらの掛け算を表現する．
たとえば，それぞれのベクトルに異なるシンボルが付加されている場合はテンソル積に，それぞれのベクトルの同じ位置に同じシンボルが付加されている場合はアダマール積に，それぞれのベクトルの上下に同じシンボルが付加されている場合は内積になる．
添字には上添字と下添字の二種類がある．
この位置関係によってアダマール積となるか内積となるか区別する．
このようにシンボリックな添字を使って，テンソルのかけあわせ方を指定する記法は，添字記法と呼ばれている．

$
\begin{pmatrix}
a_{1} & a_{2}
\end{pmatrix}_{i}
\begin{pmatrix}
b_{1} & b_{2}
\end{pmatrix}_{j}
=
\begin{pmatrix}
a_{1} b_{1} & a_{1} b_{2} \\
a_{2} b_{1} & a_{2} b_{2}
\end{pmatrix}_{ij}
$

$
\begin{pmatrix}
a_{1} & a_{2}
\end{pmatrix}_{i}
\begin{pmatrix}
b_{1} & b_{2}
\end{pmatrix}_{i}
=
\begin{pmatrix}
a_{1} b_{1} & a_{2} b_{2}
\end{pmatrix}_{i}
$

$
\begin{pmatrix}
a_{1} & a_{2}
\end{pmatrix}^{i}
\begin{pmatrix}
b_{1} & b_{2}
\end{pmatrix}_{i}
=
a_{1} b_{1} + a_{2} b_{2}
$

行列のかけ算も添字記法をつかってうまく表現できる．
内積をとるAの列とBの行について，それぞれ上下を逆にし，同じシンボルにすれば，行列の掛け算は表現できる．

$
\begin{pmatrix}
a_{11} & a_{12} \\
a_{21} & a_{22}
\end{pmatrix}^{i}_{\ j}
\begin{pmatrix}
b_{11} & b_{12} \\
b_{21} & b_{22}
\end{pmatrix}^{j}_{\ k}
=
\begin{pmatrix}
a_{11} b_{11} + a_{12} b_{21} & a_{11} b_{12} + a_{12} b_{22} \\
a_{21} b_{11} + a_{22} b_{21} & a_{21} b_{12} + a_{22} b_{22}
\end{pmatrix}^{i}_{\ k}
$

\section{添字記法を導入する手法}

以下の3つのアイデアが組み合わせることにより，添字記法が導入できる．
\begin{enumerate}
\item スカラー関数とテンソル関数
\item スカラー関数を適用したときの添字の簡約規則
\item テンソル関数には引数のテンソルの添字の情報も渡される
\end{enumerate}

\subsection{スカラー関数とテンソル関数}

スカラー関数は，スカラー値について定義された関数という意味で，テンソルを引数にとった場合，テンソルの成分ごとに処理をマップする．
テンソル関数は，テンソル値について定義された関数という意味で，テンソルを引数にとった場合，テンソルをそのまま引数としてとる．
スカラー関数は，仮引数の先頭にドルマークをつける．
テンソル関数は，仮引数の先頭にパーセントをつける．
ここでは，2引数をとってそのタプルを返す関数により，スカラー関数とテンソル関数をデモしている．

\begin{lstlisting}[language=egison,numbers=none]
(\$u $v -> (u, v)) [| a, b |]_i [| x, y |]_j
-- [| [| (a, x), (a, y) |],
        [| (b, x), (b, y) |] |]_i_j

(\%u %v -> (x, y)) [| a, b |]_i [| x, y |]_j
-- ([| a, b |]_i, [| x, y |]_j)
\end{lstlisting}

\begin{figure}[t]
 \begin{subfigure}[b]{1.0\linewidth}
  {\footnotesize
\verb|(define $min (lambda [$x $y] (if (less-than? x y) x y)))|
  }
  \caption{\texttt{min}関数の定義}
  \label{fig:minDef}
 \end{subfigure}
  \medskip
 \begin{subfigure}[b]{1.0\linewidth}
  {\small
    $min(\begin{pmatrix} 1 \\ 2 \\ 3 \\ \end{pmatrix}_{i},  \begin{pmatrix} 10 \\ 20 \\ 30 \\ \end{pmatrix}_{j}) = \begin{pmatrix} min(1,10) & min(1,20) & min(1,30) \\ min(2,10) & min(2,20) & min(2,30) \\ min(3,10) & min(3,20) & min(3,30) \\ \end{pmatrix}_{ij}
    = \begin{pmatrix} 1 & 1 & 1 \\ 2 & 2 & 2 \\ 3 & 3 & 3 \\ \end{pmatrix}_{ij}$
    }
  \caption{異なる添字をもつベクトルにたいする\texttt{min}関数の適用}
  \label{fig:minDiff}
 \end{subfigure}
  \medskip
 \begin{subfigure}[b]{1.0\linewidth}
  {\small
    $min(\begin{pmatrix} 1 \\ 2 \\ 3 \\ \end{pmatrix}_{i},  \begin{pmatrix} 10 \\ 20 \\ 30 \\ \end{pmatrix}_{i}) = \begin{pmatrix} min(1,10) & min(1,20) & min(1,30) \\ min(2,10) & min(2,20) & min(2,30) \\ min(3,10) & min(3,20) & min(3,30) \\ \end{pmatrix}_{ii}
  = \begin{pmatrix} min(1,10) \\ min(2,20) \\ min(3,30) \end{pmatrix}_{i} = \begin{pmatrix} 1 \\ 2 \\ 3 \end{pmatrix}_{i}$ 
    }
  \caption{同じ添字をもつベクトルにたいする\texttt{min}関数の適用}
  \label{fig:minSame}
 \end{subfigure}
 \caption{\texttt{min}関数の定義と適用例}
 \label{fig:minFunction}
\end{figure}

\begin{figure}[t]
 \begin{subfigure}[b]{1.0\linewidth}
  {\footnotesize
\verb|(define $. (lambda [%t1 %t2] (contract + (* t1 t2))))|
  }
  \caption{``\texttt{.}''関数の定義}
  \label{fig:dotDef}
 \end{subfigure}
  \medskip
 \begin{subfigure}[b]{1.0\linewidth}
  {\small
  $\begin{pmatrix} 1 \\ 2 \\ 3 \\ \end{pmatrix}^{i} \cdot \begin{pmatrix} 10 \\ 20 \\ 30 \\ \end{pmatrix}_{i} = contract(+, \begin{pmatrix} 10 & 20 & 30 \\ 20 & 40 & 60 \\ 30 & 60 & 90 \\ \end{pmatrix}^{i}_{\;i})
  = 10 + 40 + 90 = 140$
  \\ \medskip \\
  $\begin{pmatrix} 1 \\ 2 \\ 3 \\ \end{pmatrix}_{i} \cdot \begin{pmatrix} 10 \\ 20 \\ 30 \\ \end{pmatrix}_{i} = contract(+, \begin{pmatrix} 10 \\ 40 \\ 90 \\ \end{pmatrix}_{i}) = \begin{pmatrix} 10 \\ 40 \\ 90 \\ \end{pmatrix}_{i}$
  \\ \medskip \\
  $\begin{pmatrix} 1 \\ 2 \\ 3 \\ \end{pmatrix}_{i} \cdot \begin{pmatrix} 10 \\ 20 \\ 30 \\ \end{pmatrix}_{j} = contract(+, \begin{pmatrix} 10 & 20 & 30 \\ 20 & 40 & 60 \\ 30 & 60 & 90 \\ \end{pmatrix}_{ij})
  = \begin{pmatrix} 10 & 20 & 30 \\ 20 & 40 & 60 \\ 30 & 60 & 90 \\ \end{pmatrix}_{ij}$ \\
  }
  \caption{``\texttt{.}''関数の適用}
  \label{fig:dotApp}
 \end{subfigure}
 \caption{``\texttt{.}''関数の定義と適用例}
 \label{fig:dotFunction}
\end{figure}

テンソルの添字記法のプログラミングへの導入は，(i)\emph{スカラー仮引数}と\emph{テンソル仮引数}という2種類の仮引数の概念の導入と，(ii)適切な添字の簡約ルールによってなされる．\cite{egi2017scalar}

まず，スカラー仮引数とテンソル仮引数という2種類の仮引数の概念を説明する．
スカラー仮引数とテンソル仮引数は，\emph{スカラー関数}と\emph{テンソル関数}という2種類の関数を定義するためにそれぞれ使われる．

スカラー関数は，テンソルを引数にとったとき，成分ごとに処理をマップする関数のことをいう．
たとえば，``\verb|+|''や，``\verb|-|''，``\verb|*|''，``\verb|/|''，``\verb|∂/∂|''はスカラー関数である．
図\ref{fig:minFunction}は，スカラー関数の例として，\texttt{min}関数を定義とその適用例を紹介している．
Egisonでは、Schemeと同様\texttt{lambda}式を用いて関数を定義する．
``\verb|$|''が先頭に付加されている仮引数はスカラー仮引数となり，テンソルが引数として渡された場合，自動で成分ごとに関数がマップされる．
成分ごとに関数の処理がマップされた結果，図\ref{fig:minDiff}のように関数はテンソル積と同じ形で適用される．
テンソルの添字の簡約ルールはシンプルである．
図\ref{fig:minSame}のようにスカラー関数を適用した結果，同じシンボルの添字が2つ以上付加されたテンソルが生成された場合，その対角成分からなるテンソルに簡約されるというルールだけである．

テンソル関数は，テンソルを引数にとったとき，テンソルをテンソルとしてそのまま処理する関数のことをいう．
たとえば，行列式を計算する関数や，テンソル同士のかけ算のための関数は，テンソルの成分ごとに自動でマップされては記述できないため，テンソル関数として定義する必要がある．
図\ref{fig:dotFunction}は，テンソル関数の例として，テンソルのかけ算をする``\texttt{.}''関数の定義と適用例を紹介している．
``\verb|%|''が先頭に付加されている仮引数はテンソル仮引数となり，テンソルが引数として渡された場合，テンソルとしてそのまま渡される．
図\ref{fig:dotApp}のように，Egisonは上添字と下添字の両方を組み込みでサポートしている．
上下添字で同じシンボルが使われていた場合は，組み込み関数\texttt{contract}を使って縮約することができる．

上添字は``\verb|~|''，下添字は``\verb|_|''を使ってプログラム上で表現される．
たとえば，図\ref{fig:dotApp}の一つ目の計算例は以下のようにインタプリタ上で実行できる．

{\footnotesize
\begin{verbatim}
> (. [| 1 2 3 |]~i [| 10 20 30 |]_i)
140
\end{verbatim}
}

\subsection{省略された添字の補完}\label{completion}

添字が省略された場合の添字補完のルールを適切に設定すると，微分形式のための演算子を簡潔に定義できる．
本節では，添字補完のルールだけを説明する．
微分形式の演算子の定義の具体例は，\ref{differential-form}節で紹介する．

以下，\texttt{A}，\texttt{B}を$2$階のテンソル，つまり行列であるとする．
添字を省略した状態でスカラー関数に\texttt{A}，\texttt{B}を同時に適用すると，下記左のように同じ組み合わせの添字が補完される．
関数適用の前に``\verb|!|''が付加されていた場合は，下記右のように違う添字が補完される．

\begin{multicols}{2}

{\footnotesize
\begin{verbatim}
(+ A B) ; (+ A_t1_t2 B_t1_t2)
\end{verbatim}
}

\columnbreak

{\footnotesize
\begin{verbatim}
!(+ A B) ; (+ A_t1_t2 B_t3_t4)
\end{verbatim}
}

\end{multicols}

あとで\ref{differential-form}節で紹介するように，適切にプログラミングすれば，``\verb|!|''が使われるのは，微分形式の演算子の定義の中だけになる．
そのため，基本的には，ユーザーは上記左のルールだけを意識すればよい．


\section{\texttt{generateTensor}式によるテンソルの生成}

\lstinline{generateTensor}式は第一引数にリストを引数にとる関数，第二引数に生成したいテンソルのサイズをとる．
たとえば，すべての成分が$1$であるサイズが$(3, 4)$の行列を生成するには以下のように\lstinline{generateTensor}式を書く．
\begin{lstlisting}[language=egison,numbers=none]
generateTensor [2]#1 [3, 4]
-- [| [| 1, 1, 1, 1 |], [| 1, 1, 1, 1 |], [| 1, 1, 1, 1 |] |]
\end{lstlisting}

無名\lstinline{match}関数と組み合わせると短い記述でいろいろなテンソルを生成できる．
たとえば単位行列を生成するには以下のようにすればよい．
\begin{lstlisting}[language=egison,numbers=none]
generateTensor [2]#(\match as list integer with
                      | [$n, #n] -> 1
                      | _ -> 0)
               [3, 3]
-- [| [| 1, 0, 0 |], [| 0, 1, 0 |], [| 0, 0, 1 |] |]
\end{lstlisting}


\section{テンソルの対称性・歪対称性の宣言}

\subsection{動機}

テンソルが対称性をもつとは，下記の例のように，上三角の成分と下三角の成分が一致することをいう．
$
\begin{pmatrix}
a_{11} & a_{12} & a_{13} \\
a_{12} & a_{22} & a_{23} \\
a_{13} & a_{23} & a_{33} \\
\end{pmatrix}
$

テンソルが歪対称性をもつとは，下記の例のように，上三角の成分と下三角の成分の符号が反転することをいう．
テンソルが歪対称性をもつとき，対角成分は0となる．
$
\begin{pmatrix}
0 & a_{12} & a_{13} \\
-a_{12} & 0 & a_{23} \\
-a_{13} & -a_{23} & 0 \\
\end{pmatrix}
$

数学や物理で現れるテンソルの多くは対称性をもつ．
たとえばリーマン曲率テンソル$R_{abcd}$は$R_{abcd} = -R_{abdc}$，$R_{abcd} = -R_{bacd}$，$R_{abcd} = R_{cdab}$という対称性をもつ．
数学者や物理学者はこれらの対称性を意識しながら適宜計算を省略する．
コンピュータ・プログラムでも対称性による計算の省略をエレガントに表現したい．

$
R_{abcd} = g_{ai} \left( \frac{\partial \Gamma^{i}_{\ bd}}{\partial x_{c}} -  \frac{\partial \Gamma^{i}_{\ bc}}{\partial x_{d}} + \Gamma^{j}_{\ bd} \Gamma^{i}_{\ jc} - \Gamma^{j}_{\ bc} \Gamma^{i}_{\ jd} \right)
$


\subsection{設計}

リーマン曲率テンソルには，上記の対称性に加えて第一ビアンキ恒等式$R_{abcd} + R_{acdb} + R_{adbc} = 0$とよばれる対称性もある．
この第一ビアンキ恒等式の表現として，$R_{a[bcd]} = 0$という表記がある．
この表記を拡張してプログラムに導入するというアイデアを得た．

具体例をみせながら，その表記を紹介する．
リーマン曲率テンソル$R_{abcd}$は以下のように定義できる．
\begin{lstlisting}[language=egison,numbers=none]
def R_a_b_c_d := withSymbols [i, j]
  g_a_i . (`$\partial$`/`$\partial$` `$\Gamma$`~i_b_d x~c - `$\partial$`/`$\partial$` `$\Gamma$`~i_b_c x~d +
           `$\Gamma$`~j_b_d . `$\Gamma$`~i_j_c - `$\Gamma$`~j_b_c . `$\Gamma$`~i_j_d)
\end{lstlisting}
$R_{abcd} = -R_{abdc}$という歪対称性を宣言するには，以下のように添字\lstinline{_a_b}を角括弧\lstinline{[]}で囲む．
\begin{lstlisting}[language=egison,numbers=none]
def R[_a_b]_c_d := withSymbols [i, j]
  g_a_i . (`$\partial$`/`$\partial$` `$\Gamma$`~i_b_d x~c - `$\partial$`/`$\partial$` `$\Gamma$`~i_b_c x~d +
           `$\Gamma$`~j_b_d . `$\Gamma$`~i_j_c - `$\Gamma$`~j_b_c . `$\Gamma$`~i_j_d)
\end{lstlisting}
上記のように$R_{abcd}$を定義すると，たとえば$R_{21cd}$を参照したとき，$R_{12cd}$の計算結果の符号を反転した値を返すようになる．
対称性を複数同時に宣言することもできる．
$R_{abcd} = -R_{abdc}$と$R_{abcd} = -R_{bacd}$を同時に宣言するには，以下のように添字\lstinline{_a_b}と\lstinline{_c_d}の両方を角括弧\lstinline{[]}で囲む．
\begin{lstlisting}[language=egison,numbers=none]
def R[_a_b][_c_d] := withSymbols [i, j]
  g_a_i . (`$\partial$`/`$\partial$` `$\Gamma$`~i_b_d x~c - `$\partial$`/`$\partial$` `$\Gamma$`~i_b_c x~d +
           `$\Gamma$`~j_b_d . `$\Gamma$`~i_j_c - `$\Gamma$`~j_b_c . `$\Gamma$`~i_j_d)
\end{lstlisting}
こうすると，たとえば$R_{2121}$を参照したとき，$R_{1212}$の計算結果を返すようになる．

$R_{abcd} = R_{cdab}$のような複雑な対称性も記述できる．
この対称性を記述するためには添字\lstinline{_a_b}と\lstinline{_c_d}をそれぞれひとまとまりで扱うために丸括弧\lstinline{()}で囲む．
そして，これらを波括弧\lstinline|{}|で囲む．
\begin{lstlisting}[language=egison,numbers=none]
def R{(_a_b)(_c_d)} := withSymbols [i, j]
  g_a_i . (`$\partial$`/`$\partial$` `$\Gamma$`~i_b_d x~c - `$\partial$`/`$\partial$` `$\Gamma$`~i_b_c x~d +
           `$\Gamma$`~j_b_d . `$\Gamma$`~i_j_c - `$\Gamma$`~j_b_c . `$\Gamma$`~i_j_d)
\end{lstlisting}

対称性の宣言には三種類の括弧を使い分ける．
波括弧\lstinline|{}|は対称性を宣言するために使う．
角括弧\lstinline|[]|は歪対称性を宣言するために使う．
丸括弧\lstinline|()|は複数の添字をひとまとまりにするために使う．

三つの対称性$R_{abcd} = -R_{abdc}$，$R_{abcd} = -R_{bacd}$，$R_{abcd} = R_{cdab}$を同時に宣言することもできる．
\begin{lstlisting}[language=egison,numbers=none]
def R{[_a_b][_c_d]} := withSymbols [i, j]
  g_a_i . (`$\partial$`/`$\partial$` `$\Gamma$`~i_b_d x~c - `$\partial$`/`$\partial$` `$\Gamma$`~i_b_c x~d +
           `$\Gamma$`~j_b_d . `$\Gamma$`~i_j_c - `$\Gamma$`~j_b_c . `$\Gamma$`~i_j_d)
\end{lstlisting}

この表記法の良いところは，\lstinline{:=}の右辺のプログラムを全く変更せずに対称性を宣言できるところである．
%Wolfram言語(Mathematica)にもテンソルの対称テンソルを定義するための機能があるのですが，対称性を考慮したデータ構造をもつテンソルを生成するために右辺のプログラムは複雑化してしまう．

\subsection{さらに具体例}

さきほど，この手法の利点として，\lstinline{:=}の右辺のプログラムを全く変更する必要がないと述べたが，\lstinline{:=}の右辺のプログラムが簡略化できるケースもある．
たとえば，下記のようにテンソルの各ブロックごとに異なる式を使って値が定義されるテンソルがそのようなケースに該当する．
下記のプログラムは\ref{math}節の共同研究のためのプログラム\footnote{\url{https://github.com/egisatoshi/EMR-Paper-Computation/blob/master/program/thurston.egi}}からとってきた．
\begin{lstlisting}[language=egison,numbers=none]
def R'_i_j_k_l :=
  generateTensor
    (\match as list integer with
       | [#1, #1, _, _] -> 0
       | [_, _, #1, #1] -> 0
       | [#1, $b, #1, $d] -> -1 * p^2 * `$\delta$`~(b - 1)_(d - 1)
       | [$a, #1, #1, $d] ->      p^2 * `$\delta$`~(a - 1)_(d - 1)
       | [#1, $b, $c, #1] ->      p^2 * g_(b - 1)_(c - 1)
       | [$a, #1, $c, #1] -> -1 * p^2 * g_(a - 1)_(c - 1)
       | [#1, $b, $c, $d] -> -1 * p * `$\nabla$`J_(b - 1)_(c - 1)~(d - 1)
       | [$a, #1, $c, $d] ->      p * `$\nabla$`J_(a - 1)_(c - 1)~(d - 1)
       | [$a, $b, #1, $d] -> -1 * p * `$\nabla$`J~(d - 1)_(a - 1)_(b - 1)
       | [$a, $b, $c, #1] ->      p * `$\nabla$`J_(c - 1)_(a - 1)_(b - 1)
       | [$a, $b, $c, $d] -> R_(a - 1)_(b - 1)_(c - 1)~(d - 1)
                             + -1 * p^2 * J_(b - 1)_(c - 1) * J_(a - 1)~(d - 1)
                             +      p^2 * J_(a - 1)_(c - 1) * J_(b - 1)~(d - 1)
                             +  2 * p^2 * J_(a - 1)_(b - 1) * J_(c - 1)~(d - 1))
    [5, 5, 5, 5]
\end{lstlisting}
対称性により条件分岐を大きく減らすことができる．
\begin{lstlisting}[language=egison,numbers=none]
def R'{[_i_j][_k_l]} :=
  generateTensor
    (\match as list integer with
       | [#1, #1, _, _] -> 0
       | [#1, $b, #1, $d] -> -1 * p^2 * `$\delta$`~(b - 1)_(d - 1)
       | [#1, $b, $c, $d] -> -1 * p * `$\nabla$`J_(b - 1)_(c - 1)~(d - 1)
       | [$a, $b, $c, $d] -> R_(a - 1)_(b - 1)_(c - 1)~(d - 1)
                             + -1 * p^2 * J_(b - 1)_(c - 1) * J_(a - 1)~(d - 1)
                             +      p^2 * J_(a - 1)_(c - 1) * J_(b - 1)~(d - 1)
                             +  2 * p^2 * J_(a - 1)_(b - 1) * J_(c - 1)~(d - 1))
    [5, 5, 5, 5]
\end{lstlisting}


\subsection{テンソル関連の内部表現}

\lstinline{Symbol}データコンストラクタの第二引数

\lstinline{symbol}パターンコンストラクタの第二引数
